% GENERATED FILE. Edit canonical lessons, then run npm run generate:books.
% canonical-source-hash: fnv1a64:64fd8f20bf5beb68
% canonical-lessons: GE-C492-verlieben, GE-C492-flirten, GE-C492-kuscheln, GE-C492-kuessen, GE-C492-kuss

\chapter{Verliebt}
\label{ch:ge-verliebt}
\hlchaptermodality{492}

\begin{chapteropening}
\textbf{By the end of this chapter:} \emph{I can talk about falling in love.}

Complete the last lesson of chapter 492: I can talk about falling in love.
\end{chapteropening}

\section[verlieben]{verlieben — to fall in love (sich verlieben)}

\label{lesson:GE-C492-verlieben}



\begin{quote}
Before the new one: say the German for spouse, then the German for offspring.
\end{quote}

\subsection*{You'll want to know: verlieben}
\textbf{verlieben} — \textquotedblleft{}to fall in love (sich verlieben)\textquotedblright{}.

\begin{quote}
Im Sommer hat sie sich in ihren Tanzlehrer verliebt.
\end{quote}

\begin{quote}
\emph{In the summer she fell in love with her dance teacher.}
\end{quote}

Reflexive: \emph{sich verlieben in} + accusative. \emph{Verliebt sein} — to be in love.

\begin{cousinweb}
\emph{verlieben} is built on \emph{ver-} and \emph{lieben} \textquotedblleft{}to love\textquotedblright{} (EtymDB lists both), a Germanic verb related to \emph{lieb} \textquotedblleft{}dear\textquotedblright{}, from Proto-Germanic \emph{leubaz}.
\end{cousinweb}

\begin{grammarlens}[title={What you've built}]
One more word for telling, explaining and describing at B1.
\end{grammarlens}

\subsection*{Your turn}
\begin{itemize}
\raggedright
  \item \emph{Say it:} \emph{verlieben}
  \item \emph{Say it:} \emph{verlieben}, once more
  \item \emph{Say it:} a sentence of your own with \emph{verlieben}
  \item \emph{Recall:} say the German for spouse, then the German for offspring, then say \emph{verlieben} again
\end{itemize}

\begin{tcolorbox}[breakable,colback=teal!4,colframe=teal!35!black,title={Before you move on}]
What does verlieben mean? (\textquotedblleft{}To fall in love (sich verlieben)\textquotedblright{}.) Say it once more.
\end{tcolorbox}
\section[flirten]{flirten — to flirt}

\label{lesson:GE-C492-flirten}



\begin{quote}
Before the new one: say the German for offspring, then the German for to fall in love.
\end{quote}

\subsection*{You'll want to know: flirten}
\textbf{flirten} — \textquotedblleft{}to flirt\textquotedblright{}.

\begin{quote}
Auf der Party flirtete er mit allen.
\end{quote}

\begin{quote}
\emph{At the party he flirted with everybody.}
\end{quote}

Regular: \emph{hat geflirtet}.

\begin{cousinweb}
\emph{flirten} was borrowed from English \emph{to flirt}; EtymDB has no entry for the German verb.
\end{cousinweb}

\begin{grammarlens}[title={What you've built}]
One more word for telling, explaining and describing at B1.
\end{grammarlens}

\subsection*{Your turn}
\begin{itemize}
\raggedright
  \item \emph{Say it:} \emph{flirten}
  \item \emph{Say it:} \emph{flirten}, once more
  \item \emph{Say it:} a sentence of your own with \emph{flirten}
  \item \emph{Recall:} say the German for offspring, then the German for to fall in love, then say \emph{flirten} again
\end{itemize}

\begin{tcolorbox}[breakable,colback=teal!4,colframe=teal!35!black,title={Before you move on}]
What does flirten mean? (\textquotedblleft{}To flirt\textquotedblright{}.) Say it once more.
\end{tcolorbox}
\section[kuscheln]{kuscheln — to cuddle}

\label{lesson:GE-C492-kuscheln}



\begin{quote}
Before the new one: say the German for to fall in love, then the German for to flirt.
\end{quote}

\subsection*{You'll want to know: kuscheln}
\textbf{kuscheln} — \textquotedblleft{}to cuddle\textquotedblright{}.

\begin{quote}
Am Abend kuscheln wir auf dem Sofa.
\end{quote}

\begin{quote}
\emph{In the evening we cuddle on the sofa.}
\end{quote}

\emph{Sich an jemanden kuscheln} — to snuggle up to someone.

\begin{cousinweb}
The origin of \emph{kuscheln} is not settled; EtymDB has no entry for it.
\end{cousinweb}

\begin{grammarlens}[title={What you've built}]
One more word for telling, explaining and describing at B1.
\end{grammarlens}

\subsection*{Your turn}
\begin{itemize}
\raggedright
  \item \emph{Say it:} \emph{kuscheln}
  \item \emph{Say it:} \emph{kuscheln}, once more
  \item \emph{Say it:} a sentence of your own with \emph{kuscheln}
  \item \emph{Recall:} say the German for to fall in love, then the German for to flirt, then say \emph{kuscheln} again
\end{itemize}

\begin{tcolorbox}[breakable,colback=teal!4,colframe=teal!35!black,title={Before you move on}]
What does kuscheln mean? (\textquotedblleft{}To cuddle\textquotedblright{}.) Say it once more.
\end{tcolorbox}
\section[küssen]{küssen — to kiss}

\label{lesson:GE-C492-kuessen}



\begin{quote}
Before the new one: say the German for to flirt, then the German for to cuddle.
\end{quote}

\subsection*{You'll want to know: küssen}
\textbf{küssen} — \textquotedblleft{}to kiss\textquotedblright{}.

\begin{quote}
Zum Abschied küsste er sie.
\end{quote}

\begin{quote}
\emph{He kissed her goodbye.}
\end{quote}

\emph{Sich küssen} — to kiss each other.

\begin{cousinweb}
\emph{küssen} is inherited from Old High German \emph{kussen}, from Proto-Germanic \emph{kussijaną} \textquotedblleft{}to kiss\textquotedblright{} (EtymDB); English \emph{kiss} is a cousin.
\end{cousinweb}

\begin{grammarlens}[title={What you've built}]
One more word for telling, explaining and describing at B1.
\end{grammarlens}

\subsection*{Your turn}
\begin{itemize}
\raggedright
  \item \emph{Say it:} \emph{küssen}
  \item \emph{Say it:} \emph{küssen}, once more
  \item \emph{Say it:} a sentence of your own with \emph{küssen}
  \item \emph{Recall:} say the German for to flirt, then the German for to cuddle, then say \emph{küssen} again
\end{itemize}

\begin{tcolorbox}[breakable,colback=teal!4,colframe=teal!35!black,title={Before you move on}]
What does küssen mean? (\textquotedblleft{}To kiss\textquotedblright{}.) Say it once more.
\end{tcolorbox}
\section[der Kuss]{der Kuss — kiss}

\label{lesson:GE-C492-kuss}



\begin{quote}
Before the new one: say the German for to cuddle, then the German for to kiss.
\end{quote}

\subsection*{You'll want to know: der Kuss}
\textbf{der Kuss} — \textquotedblleft{}kiss\textquotedblright{}.

\begin{quote}
Die Mutter gibt dem Kind einen Kuss auf die Wange.
\end{quote}

\begin{quote}
\emph{The mother gives the child a kiss on the cheek.}
\end{quote}

Plural: \emph{die Küsse}.

\begin{cousinweb}
\emph{Kuss} is inherited from Middle High German \emph{kus} (EtymDB), from the same Germanic source as \emph{küssen}, Proto-Germanic \emph{kussijaną}.
\end{cousinweb}

\begin{grammarlens}[title={What you've built}]
One more word for telling, explaining and describing at B1.
\end{grammarlens}

\subsection*{Your turn}
\begin{itemize}
\raggedright
  \item \emph{Say it:} \emph{der Kuss}
  \item \emph{Say it:} \emph{der Kuss}, once more
  \item \emph{Say it:} a sentence of your own with \emph{der Kuss}
  \item \emph{Recall:} say the German for to cuddle, then the German for to kiss, then say \emph{der Kuss} again
\end{itemize}

\begin{tcolorbox}[breakable,colback=teal!4,colframe=teal!35!black,title={Before you move on}]
What does der Kuss mean? (\textquotedblleft{}Kiss\textquotedblright{}.) Say it once more.
\end{tcolorbox}
